\section{Modulus-Dimension Switching}
\label{sec:modexpansion}



The main results of this section are Corollaries~\ref{cor:mod-reduction} and~\ref{cor:mod-dim-tradeoff} below.
Both are special cases of the following technical theorem.  
%\dnote{I added this and moved from~$B$-bounded to~$(B,\delta)$-bounded} 
We
say that a distribution~$\calD$ over~$\Z^{n}$ is $(B,\delta)$-bounded
for some reals~$B,\delta\geq 0$ if the probability that~$\vec{x} \gets
\calD$ has norm greater than $B$ is at most $\delta$. 

\begin{theorem}
  \label{thm:mod-dim-switch}
  Let $m,n, n',q,q' \geq 1$ be integers, let
  $\matG \in \Z^{n' \times n}$ be such that the lattice $\Lambda =
  \frac{1}{q'} \matG^{T} \Z^{n'} + \Z^{n}$ has a known basis $\matB$,
  and let $\calD$ be an arbitrary $(B,\delta)$-bounded distribution over
  $\Z^{n}$.  Let $\alpha, \beta > 0$ and~$\eps \in (0,1/2)$ satisfy 
	\[\beta^{2} \geq
  \alpha^{2} + (4/\pi) \ln(2n(1+1/\eps)) \cdot (\max\set{q^{-1}, \length{\gs{\matB}}} \cdot B)^{2}.\]  
	Then there is an efficient (transformation) reduction
  from~$\lwe_{n,m,q,\le \alpha}(\calD)$ to $\lwe_{n',m,q',\leq \beta}(\matG
  \cdot \calD)$ that reduces the advantage by at most $\delta + 14 \eps m$.
\end{theorem}

Here we use the notation
$\|\widetilde \matB\|$ from Lemma~\ref{lem:gpv}. We also note that if needed, the distribution on secrets produced
by the reduction can always be turned into the uniform distribution on $\Z_{q'}^{n'}$, as mentioned after Definition~\ref{def:lwe}.
Also, we recall that there exists an elementary reduction from $\lwe_{n',q',\leq \beta}$ 
to~$\lwe_{n',q',\beta}$ (see Lemma~\ref{lem:lelwe}).

Here we state two important corollaries of the theorem.  The first
corresponds to just modulus reduction (the \lwe dimension is
preserved), and is obtained by letting $n'=n$, $\matG=\matI$ be the
$n$-dimensional identity matrix, and $\matB = \matI/q'$.  
For example, we can take $q \ge q' \ge \sqrt{2\ln(2n(1+1/\eps))} \cdot
(B/\alpha)$ and $\beta=\sqrt{2} \alpha$, which corresponds to reducing
an arbitrary modulus to almost~$B/\alpha$, while increasing the
initial error rate~$\alpha$ by just a small constant factor.

\begin{corollary}
  \label{cor:mod-reduction}
  For any $m,n \geq 1$, $q \geq q' \geq 1$, $(B,\delta)$-bounded distribution~$\calD$ over~$\Z^{n}$, $\alpha,\beta>0$ and~$\eps \in (0,1/2)$ such that
  \[
	\beta^{2} \geq \alpha^{2} + (4/\pi) \ln(2n(1+1/\eps))
  \cdot (B/q')^2,
	\]
	there is an efficient reduction from
  $\lwe_{n,m,q,\le \alpha}(\calD)$ to 
	$\lwe_{n,m,q',\le \beta}(\calD)$ that reduces the advantage by at most $\delta + 14 \eps m$.
\end{corollary}

In particular, by using the normal form of \lwe (Lemma~\ref{lem:normalform}), in which the secret has distribution $\calD = D_{\Z^{n}, \sqrt{2} \alpha q}$, we can switch to a power-of-2 modulus with only a small loss in the noise rate, as described in the following
corollary. Together with the known search-to-decision reduction (Theorem~\ref{thm:searchtodecisionmicpei}), this extends the known hardness
of (decision) $\lwe$ to \emph{any} modulus $q$.
Here we use that $\calD = D_{\Z^{n}, r}$
is $(C r \sqrt{n \log(n/\delta)}, \delta)$-bounded for some universal constant $C>0$, which
follows by taking union bound over the $n$ coordinates.
(Alternatively, one could use that it is~$(r\sqrt{n},2^{-n})$-bounded, as follows
from~\cite[Lemma~1.5]{banaszczyk93:_new}, leading to a slightly tighter statement for large $n$.)

\begin{corollary}
  \label{cor:mod-reduction-2}
  Let $\delta \in (0,1/2)$, $m \ge n \geq 1$, $q' \ge 25$. Let also $q \in [q',2q')$ be the smallest power of $2$ not smaller than $q'$
	and $\alpha \ge \sqrt{\ln(2n(1+16/\delta)/\pi)}/q$.
	There exists an efficient (transformation) reduction from $\lwe_{n,m,q, \alpha}$ to $\lwe_{n,m',q', \le \beta}$
	where $m'=m-(16 n + 4 \ln\ln q)$ and 
	\[
	\beta = C \alpha \sqrt{n} \sqrt{\log(n/\delta)\log(m/\delta)}
	\]
	for some universal constant $C>0$, that turns advantage of $\zeta$ into an advantage of at least $(\zeta-\delta)/4$.
\end{corollary}


  Another corollary illustrates a modulus-dimension tradeoff.  Assume
  $n=kn'$ for some $k \ge 1$, and let~$q' = q^{k}$.  Let~$\vec{G} =
  \vec{I}_{n'} \otimes \vec{g}$, where~$\vec{g} = (1,q,q^2, \ldots,
  q^{k-1})^T \in \Z^k$.  We then have $\Lambda = q^{-k} \vec{G}^T
  \Z^{n'} + \Z^n$.  A basis of~$\Lambda$ is given by
\[
\vec{B}= \vec{I}_{n'} \otimes 
\begin{bmatrix}
q^{-1} & q^{-2} & \cdots & q^{-k} \\
       & q^{-1} & \cdots & q^{1-k} \\
       &        & \ddots & \vdots \\
       &		&		 & q^{-1}
     \end{bmatrix} \in \R^{n \times n};
\]
this is since the column vectors of~$\vec{B}$ belong to~$\Lambda$ and
the determinants match.  Orthogonalizing from left to right, we have
$\gs{\matB} = q^{-1} \matI$ and so $\length{\gs{\matB}}=q^{-1}$.  We
therefore obtain the following corollary, showing that we can trade
off the dimension against the modulus, holding~$n \log q = n' \log q'$
fixed.  For example, letting~$\calD = D_{\Z^{n}, \alpha q}$
(corresponding to a secret in normal form, see Lemma~\ref{lem:normalform}), which is $(\alpha q
\sqrt{n}, 2^{-n})$-bounded, 
the reduction increases the error rate by about
a~$\sqrt{n}$ factor. 

\begin{corollary}
  \label{cor:mod-dim-tradeoff}
  For any $n, m,q \ge 1$, $k \ge 1$ that divides~$n$,  $(B,\delta)$-bounded distribution~$\calD$ over~$\Z^{n}$, 
	$\alpha,\beta>0$, and~$\eps \in (0,1/2)$ such 
that
\[
\beta^{2} \geq \alpha^{2} + (4/\pi)
  \ln(2n(1+1/\eps)) \cdot (B/q)^{2},
\]
  there is an efficient reduction from~$\lwe_{n,m,q,\le \alpha}(\calD)$
  to~$\lwe_{n/k,m,q^k,\le \beta}(\vec{G} \cdot \calD)$ that reduces the advantage by at most $\delta+14 \eps m$, 
where~$\vec{G} = \vec{I}_{n/k} \otimes  (1,q,q^2, \ldots, q^{k-1})^T$.
\end{corollary}

\noindent Theorem~\ref{thm:mod-dim-switch} follows immediately from
the following lemma.


\begin{lemma}
  \label{lem:reverse-main}
  Adopt the notation of Theorem~\ref{thm:mod-dim-switch}, and let 
\[
	r
  \geq \max\set{q^{-1}, \length{\gs{\matB}}} \cdot
  \sqrt{2\ln(2n(1+1/\epsilon))/\pi}. 
	\]
There is an efficient 
  mapping from $\T_{q}^{n} \times \T$ to $\T_{q'}^{n'} \times \T$, which
  has the following properties:
  \begin{itemize}[itemsep=0pt]
  \item If the input is uniformly random, then the output is within
statistical distance $4 \epsilon$ from the uniform distribution.
  \item If the input is distributed according to
    $A_{q,\vecs,D_{\alpha}}$ for some $\vecs \in \Z^{n}$ with
    $\length{\vecs} \leq B$, then the output distribution is within
    statistical distance $10 \epsilon$ from $A_{q',\matG\vecs,D_{\alpha'}}$, where
    $(\alpha')^{2} = \alpha^{2} + r^{2}(\length{\vecs}^{2}+B^{2}) \leq
    \alpha^{2} + 2(rB)^{2}$. 
%\dnote{I think we use the condition $\length{\vecs} \leq B$ only for that last inequality. 
%If this is indeed the case, I would be in favour of removing the condition, and commenting at the end of this item that 
%if~$\length{\vecs} \leq B$, then $\alpha^{2} + r^{2}(\length{\vecs}^{2}+B^{2}) \leq
%    \alpha^{2} + 2(rB)^{2}$}\onote{don't we also need it when we apply Lemma~\ref{lem:smoothip} towards the end of the proof? (I didn't check too carefully)}
%\dnote{Why? we don't need $\vec{s}$ small too apply it, do we?}\onote{yes, I think we do; look carefully at the lemma...}
  \end{itemize}
\end{lemma}

\begin{proof}
  The main idea behind the reduction is to encode $\T_{q}^{n}$ into
  $\T_{q'}^{n'}$, so that the mod-$1$ inner products between vectors
  in $\T_{q}^{n}$ and a short vector $\vecs \in \Z^{n}$, and between
  vectors in $\T_{q'}^{n'}$ and $\matG \vecs \in \Z^{n'}$, are nearly
  equivalent. In a bit more detail, the
  reduction will map its input vector $\veca \in \T_{q}^{n}$ (from the
  given LWE-or-uniform distribution) to a vector $\veca' \in
  \T_{q'}^{n'}$, so that
  \[ \inner{\veca', \matG \vecs} = \inner{\matG^{T} \veca', \vecs}
  \approx \inner{\veca,\vecs} \bmod 1 \] for any (unknown) $\vecs \in
  \Z^{n}$.  To do this, it randomly samples $\veca'$ so that
  $\matG^{T} \veca' \approx \veca \bmod \Z^{n}$, where the
  approximation error will be a discrete Gaussian of parameter $r$.
  

  % In order to formally describe how the reduction samples $\veca'$, we
  % first define an $n$-dimensional lattice \[ \Lambda = q'^{-1} \cdot
  % \lat(\matG^{T}) + \Z^{n} = q'^{-1} \cdot \lat(\matI_{n'} \otimes
  % \vecg) + \Z^{n}, \] and claim that it has as a basis the
  % upper-triangular matrix \[ \matB =
  % \matI_{n'} \otimes
  % \begin{bmatrix}
  %   (q'_{1})^{-1} & \star & \cdots & \star \\
  %   & (q'_{2})^{-1} & \cdots & \star \\
  %   & & \ddots & \vdots \\
  %   & & & (q'_{k})^{-1}
  % \end{bmatrix} \in \R^{n \times n}, \] where the rightmost column of
  % the above displayed $k$-by-$k$ matrix is $q'^{-1} \vecg$, and the
  % $j$th column is obtained by multiplying the $(j+1)$st column by
  % $q'_{j+1}$ and reducing modulo $1$, for $j=k-1,\ldots, 1$.  Indeed,
  % $\matB$ is a basis of $\Lambda$ because all of its columns are in
  % $\Lambda$ by construction, and because
  % $\det(\matB)=q'^{-n'}=\abs{\Lambda/\Z^{n}}^{-n'}=\det(\Lambda)$.  Also
  % notice that because $\matB$ is upper triangular, its Gram-Schmidt
  % orthogonalization (from left to right) is just the diagonal matrix
  % with $(q'_{i})^{-1}$ in the $i$th diagonal, so $\max_{i}
  % \length{\gs{\vecb_{i}}} \leq q_{\min}^{-1}$ and $r \geq
  % \sqrt{2}\smootheps(\Lambda)$ by~\cite[Theorem
  % 3.1]{DBLP:conf/stoc/GentryPV08}\cnote{Put this in prelims...}.
  
  We can now formally define the reduction, which works as follows.
  On an input pair $(\veca, b) \in \T_{q}^{n} \times \T$, it does the
  following:
  \begin{itemize}
  \item Choose $\vecf \gets D_{\Lambda-\veca,r}$ using Lemma~\ref{lem:gpv} with basis
    $\matB$, and let $\vecv = \veca + \vecf \in
    \Lambda/\Z^{n}$. (The coset $\Lambda-\veca$ is well defined since $\veca =
      \bar{\veca}+\Z^{n}$ is some coset of $\Z^{n} \subseteq \Lambda$.) Choose a uniformly
    random solution $\veca' \in \T_{q'}^{n'}$ to the equation
    $\matG^{T} \veca' = \vecv \bmod \Z^{n}$. This can be done by computing a basis of the solution set $\matG^{T}
      \veca' = \veczero \bmod \Z^{n}$, and adding a uniform element from that set to an arbitrary solution to the equation
    $\matG^{T} \veca' = \vecv \bmod \Z^{n}$.
  \item Choose $e' \gets D_{rB}$ and let $b' = b + e' \in \T$.  
  \item Output $(\veca',b')$.
  \end{itemize}

  We now analyze the reduction.  First, if the distribution of the
  input is uniform, then it suffices to show that $\veca'$ is
  (nearly) uniformly random,
  because both $b$ and~$e'$ are independent of~$\veca'$, and $b \in
  \T$ is uniform.  To prove this claim, notice that it suffices to
  show that the coset $\vecv \in \Lambda/\Z^{n}$ is (nearly) uniformly
  random, because each $\vecv$ has the same number of solutions
  $\veca'$ to $\matG^{T} \veca' = \vecv \bmod \Z^{n}$.  Next, observe that for any $\bar{\veca} \in \T_q^n$ and $\bar{\vecf} \in \Lambda - \bar{\veca}$, we have by Lemma~\ref{lem:smoothing} 
	(using that $r \ge \smootheps(\Lambda)$ by Lemma~\ref{lem:boundonsmoothing}) that
	\begin{align}\label{eq:roundingonelatticeanother}
	  \Pr[\veca = \bar{\veca} \wedge \vecf = \bar{\vecf}] &= q^{-n} \cdot \rho_r(\bar{\vecf}) / \rho_r(\Lambda-\bar{\veca}) \nonumber \\
		& \in C \bracks[\big]{1,\tfrac{1+\epsilon}{1-\epsilon}}\cdot \rho_r(\bar \vecf).
	\end{align}
	where $C=q^{-n}/\rho_{r}(\Lambda)$ is a normalizing
  value that does not depend on $\bar{\veca}$ or $\bar{\vecf}$.
	Therefore, by summing over all $\bar \veca, \bar \vecf$ satisfying $\bar \veca + \bar \vecf = \bar \vecv$, we obtain that 
	for any $\bar{\vecv} \in \Lambda/\Z^{n}$,
  \begin{align*}
    \Pr[\vecv = \bar{\vecv}] 
    &\in C \bracks[\big]{1,\tfrac{1+\epsilon}{1-\epsilon}} \cdot
    \rho_{r}(q^{-1} \Z^{n} + \bar{\vecv}).
  \end{align*}
  Since $r
  \geq \smootheps(q^{-1} \Z^{n})$ (by
  Lemma~\ref{lem:boundonsmoothing}), Lemma~\ref{lem:smoothing} 
implies that~$\Pr[\vecv = \bar{\vecv}] \in \bracks[\big]{\tfrac{1-\epsilon}{1+\epsilon},\tfrac{1+\epsilon}{1-\epsilon}} C'$
for a constant~$C'$ that is independent of~$\bar{\vec{v}}$. By Claim~\ref{clm:separationdist}, this shows that~$\veca'$
is within statistical distance~$1-((1-\eps)/(1+\eps))^2 \le 4\epsilon$ of the uniform distribution.

  It remains to show that the reduction maps $A_{q,\vecs,D_{\alpha}}$
  to $A_{q', \matG \vecs, D_{\beta}}$.  Let the input sample from the
  former distribution be $(\veca,b=\inner{\veca,\vecs}+e)$, where $e
  \gets D_{\alpha}$.  As argued above, the output $\veca'$ is (nearly)
  uniform over~$\T_{q'}^{n'}$.  So condition now on any fixed value $\overline{\veca'} \in \T_{q'}^{n'}$ of $\veca'$,
	and let $\bar \vecv = \matG^{T} \overline{\veca'} \bmod \Z^{n}$. 	
	We have
  \[ b' = \inner{\veca,\vecs} + e + e' = \inner{\overline{\veca'}, \matG \vecs} +
  e + \inner{-\vecf, \vecs} + e' \bmod 1. \] 
	By Claim~\ref{clm:separationdist} and~\eqref{eq:roundingonelatticeanother} (and noting that if $\vecf=\bar \vecf$ then $\veca = \bar \vecv - \bar \vecf \bmod \Z^n$),
	the distribution of $-\vecf$ is within statistical distance~$1-(1-\eps)/(1+\eps) \le 2 \eps$ of
	$D_{q^{-1}\Z^{n}-\bar \vecv, r}$.  By Lemma~\ref{lem:smoothip} (using $r
  \geq \sqrt{2} \smootheps(q^{-1} \Z^{n})$ and~$\length{\vecs} \le B$), the
  distribution of $\inner{-\vecf,\vecs} + e'$ is within statistical 
distance~$6 \eps$ 
  from~$D_{t}$, where $t^{2} = r^{2}(\length{\vecs}^{2}+B^{2})$.  It
  therefore follows that $e+\inner{-\vecf,\vecs} + e'$ is within statistical distance~$6 \eps$ 
from $D_{(t^2+\alpha^2)^{1/2}}$, as required.
\end{proof}

\begin{comment}

\subsection{Application to Ring-LWE}

\cnote{I think the theorem in this section is a special case of the
  ``ring-LWE embedding'' theorem from GHPS (SCN'12), which says that
  we can embed LWE over a cyclotomic ring with index $m$ into LWE over
  a cyclotomic of index evenly divisible by $m$.  This theorem is the
  special case where $m=1$ (so the cyclotomic ring is just $\Z$).}

\begin{definition}
\label{def:rlwe}
Let~$R_{n,q}$ denote the ring~$\Z_q[x]/(x^n+1)$, with~$n$ a power
of~$2$ and~$q$ an integer. Let~$\chi$ be a distribution
over~$\mathbb{T}_n = \mathbb{T}[x]/(x^n+1)$, and~$\mathcal{D}$ be a
distribution from~$R_{n,q}$. For~$s \in R_{n,q}$, we define the
distribution~$A_{n,q,\chi,s}^R$ as follows: sample~$a \hookleftarrow
U(R_{n,q})$ and~$e \hookleftarrow \chi$; then return~$(a,
\frac{1}{q}(a\cdot s)+e) \in R_{n,q} \times \mathbb{T}_n$.


The decision \emph{Ring Learning with Errors} problem
$\rlwe_{n,q,\chi}(\mathcal{D})$ is as follows: Let~$\vec{s}$ be
sampled from~$\mathcal{D}$; the goal is to distinguish between
arbitrarily many independent samples from~$A_{n,q,\vec{s},\chi}^R$ and
the same number of independent samples from the uniform distribution
over the support of~$A^R_{n,q,\vec{s},\chi}$.
\end{definition}

The above RLWE problem is a special case of the ring learning with
errors problem that was introduced
in~\cite{DBLP:conf/eurocrypt/LyubashevskyPR10}. Since then, it has
proved useful for designing cryptographic schemes with excellent
asymptotic performance. However, all known hardness proofs for \rlwe
rely on the assumed hardness of standard worst-case lattice problems
restricted to the family of so-called ideal lattices. In our present
setup, an ideal lattice is an ideal of~$\Z[x]/(x^n+1)$.
Theorem~\ref{th:reverse} leads to a reduction from \lwe to \rlwe, and
thus provides a hardness proof for \rlwe under standard worst-case
lattice problems over all lattices.

\begin{theorem}
\label{th:rlwe}
For any~$q, \chi, \mathcal{D}$, there exists a polynomial-time
reduction from~$\lwe_{1,q,\chi}(\mathcal{D})$
to $\rlwe_{n,q,\chi'_n}(\mathcal{D}_n)$ where~$\chi'_n = \chi'+\chi' x
+ \cdots + \chi' x^{n-1}$, $\chi'$ is the sum of $n$ independent
copies of~$\chi$ and~$\mathcal{D}_n = \mathcal{D} + \mathcal{D} x +
\cdots + \mathcal{D} x^{n-1}$.
\end{theorem}

\begin{proof}
  For~$s_0,\ldots,s_{n-1} \in \Z_p$, we define the
  distribution~$A_{n,q,\chi,(s_i)_i}$: Sample~$a \hookleftarrow
  U(\Z_q)$ and~$e_0,\ldots,e_{n-1} \hookleftarrow \chi$, and
  return~$(a, \frac{1}{p}(a\cdot
  s_0)+e_0,\ldots,\frac{1}{p}(as_{n-1})+e_{n-1}) \in \Z_p \times
  \mathbb{T}^n$.  Using a standard hybrid argument, we see
  that~$\lwe_{1,q,\chi}(\mathcal{D})$ reduces to distinguishing between
  arbitrarily many independent samples from~$A_{n,q,\chi,(s_i)_i}$ and
  the same number of independent samples from the uniform distribution
  over the support of~$A_{n,q,\chi,(s_i)_i}$. We describe an efficient
  mapping of~$n$ elements~$(a_i, (b_{ij})_{0\leq j<n})_{0 \leq i<n}$
  from~$\Z_p \times \mathbb{T}^n$ to a single element in~$R_q \times
  \mathbb{T}_n$ such that the $A_{n,q,\chi,(s_i)_i}$ distribution
  (resp.\ uniform distribution over the support
  of~$A_{n,q,\chi,(s_i)_i}$) is mapped to the~$A_{n,q,\chi_n,s}^R$
  distribution (resp.\ uniform distribution over the support
  of~$A_{n,q,\chi_n,s}^R$), where~$s = s_0+s_1x + \cdots
  +s_{n-1}x^{n-1}$. The mapping is as follows:
\begin{itemize}
\item Define~$a = a_0 + a_1x + \cdots + a_{n-1}x^{n-1} \in R_{n,q}$.
\item Define~$b = \sum_{0\leq k<2n-1} (\sum_{i+j = k} b_{ij}) x^k \bmod (x^n+1)$.
\item Return~$(a,b)$.  
\end{itemize}
Clearly, the reduction runs in polynomial-time, and the uniform
distribution is mapped to the uniform distribution. Finally, in case
the input~$b_{ij}$'s are of the form~$\frac{1}{p} (a_i\cdot s_j) +e_j$
for all~$i,j$, then all the coefficients of~$b - \frac{1}{p}(a\cdot s)
\in \mathbb{T}_n$ are sums of~$n$ independent samples from~$\chi$.
\end{proof}

As a corollary of Theorems~\ref{th:reverse} and~\ref{th:rlwe}, there
exists a (classical) polynomial-time reduction from standard
worst-case problems over lattices with polynomial approximation
factors to~$\rlwe_{n,p^n,D_{\sqrt{n}\alpha}}(U(R_{n,p^n}))$.

\end{comment}

\begin{comment}
\begin{verbatim}
Also, we were wondering about a possible relationship with
   http://www.di.ens.fr/~lyubash/papers/subsetsumcrypto.pdf
(digit-wise expansion of a large integer, along columns rather than rows).

Notation: \vec* denotes an n-dimensional vector mod p (i.e., in Zp^n),
while capital letters (A,S) denote scalar values mod p^n (i.e., in
Z_{p^n}).

The basic approach is to map LWE samples (\vecr,b) in Zp^n \times Zp
to (A, b*p^{n-1} + noise) in Z_{p^n} \times Z_{p^n}, where A is chosen
from r using some *randomized* encoding (not deterministically).  Then
the amount by which b*p^{n-1} differs from the value A*S will be a
centered Gaussian, and we only need to add a small amount of extra
noise (or maybe none at all?) to smooth it out.

Details:

For A,S in Z_p^n we can write A = sum_i p^i * Ai for *integers* ai
(not necessarily in {0,...,p-1}), and similarly for S.

Then the product A*S mod p^n can be written as \veca^t \matP \vecs mod
p^n, where

\veca = (a_{n-1}, ..., a_0),   \vecs = (s0, s1, ..., s_{n-1}),  and
\matP is the fixed lower-triangular Toeplitz matrix with
\vecp=(p^{n-1}, p^{n-2}, ..., p^0) down its leftmost column.

Ideally, we would then like to choose \veca (defining A) to satisfy

  \veca^t \matP = p^{n-1} * \vecr^t    (mod p^n)

where \vecr is the "left-hand side" of the LWE challenge.

If we could do this, then the product A*S would be exactly p^{n-1} *
<r,s>.  We are given the noisy value b=<r,s>+error, so we could just
plug in p^{n-1}*b (mod p^n) to complete the reduction.  Notice that
this would keep the error rate exactly the same, and moreover the
error term would be divisible by p^{n-1}, which seems like too much to
ask.... and it is.

The problem, of course, is that it is usually not possible to solve
the above system, because \matP is highly non-invertible mod p^n.  (It
doesn't even help that the RHS is a multiple of p^{n-1}.)

Luckily, it's enough to approximately solve the system, as long as the
"error" \veca^t \matP - p^{n-1} * \vecr^t is a centered Gaussian and
\vecs is short.  And the structure of \matP makes this easy to do --
Oded's formulas are the answer.

[Chris note: the formulas come from an earlier email, with slightly
different notation.  But they are easy to work out: the entries of
\veca are chosen iteratively from discrete Gaussians over the
integers, with appropriate centers determined by the previously chosen
entries of \veca.]

Another way of going about it is to multiply both sides by \matP^-1
**over the integers**, which has a very simple form (exercise for
reader).  Then we end up with a fractional RHS which we can just
randomized-round to get an integral \veca.  Doing it this way causes
the "error" \veca^t \matP - p^{n-1} * \vecr^t to be a non-spherical
but "nearly spherical" Gaussian (because \matP is almost orthogonal),
which would suffice for our purposes.  It's not clear that this
non-adaptive solution has any important advantage here, though.

(Note: these two approaches correspond to randomized versions of
Babai's nearest-plane and "simple rounding" algorithms, respectively.)

\end{verbatim}
\end{comment}



%%% Local Variables: 
%%% mode: latex
%%% TeX-master: "lwehard"
%%% End: 
